
Holographic near-eye displays offer unique benefits to virtual and augmented reality (VR/AR) applications. For example, holographic displays can present perceptually realistic 3D images with natural parallax to the user in lightweight device form factors~\cite{maimone2017holo,Kim2022-cd, gopakumar2024ar,jang2024waveguide}. Yet, the \'etendue of holographic displays is fundamentally limited by the pixel count of the underlying spatial light modulators (SLMs), preventing current holographic near-eye displays from achieving a large field of view and eyebox simultaneously. This limitation is a fundamental barrier towards making this a practical display technology. 


Increasing the pixel count of an SLM seems like the natural solution. However, developing large-area phase-only SLMs with pixel pitches matching the small feature sizes (i.e., tens of nanometers) of analog holographic films~\cite{benton2008holographic} is simply not feasible with today's hardware solutions. To overcome this problem, \'etendue expansion techniques have been described in the literature, including those based on static, high-resolution masks~\cite{kuo2020mask, park2019photon,yu2017speckle, buckley2006mask,monin2022mask,tseng2023neural}, pupil replication~\cite{kress2020waveguide}, steered or multi-source illumination~\cite{jang2018hoe,lee2020wide,jo2022binary,lee2022envelope, monin2022tilt}, and making use of higher-diffraction orders and pupil optimization \cite{shi2024ergo, schiffers2023stochastic}. However, each of these approaches has its limitations, as mask-based systems do not have sufficient degrees of freedom to achieve a high image quality, pupil replication approaches cannot create natural 3D effects and parallax over the eyebox, and steered sources are hindered by the requirement for high-speed SLMs as well as high diffraction orders (HDOs) that fundamentally limit the image quality. As a result, none of these solutions is able to achieve high-quality 3D holography with a large \'etendue.


Our work is motivated by the hypothesis that a holographic display requires sufficient degrees of freedom to achieve a large field of view and eyebox simultaneously. In the absence of an extremely high-resolution SLM, this is only achievable using steered or multi-source illumination. We thus build on the latter approach but address its major shortcomings, HDOs and symmetric illumination copies, by introducing a dynamic, programmable amplitude modulation mechanism in the Fourier plane, after the SLM. This unique optical setup allows us to extend steered / multi-source configurations such that they modulate the frequency spectrum of the display image in a content-adaptive manner. For this purpose, we leverage a stochastic optimization approach that factors a target light field into a set of time-multiplexed phase SLM and corresponding Fourier amplitude masks that are displayed in rapid succession while being integrated by the user's eye.


Using the proposed system, we demonstrate improved 3D image quality over a large \'etendue, surpassing the performances of existing approaches in both simulation and experiment. Specifically, our contributions include 
%
\begin{itemize}
    \item A novel optical holographic display configuration that combines a time-multiplexed phase SLM near the image plane and a dynamic amplitude SLM that controls the frequency spectrum.
    \item A computer-generated holography framework that uses stochastic optimization to factor a target light field into a set of phase--amplitude image pairs.
    \item Demonstration of improved 3D image quality among high-\'etendue holographic displays.
\end{itemize}

Our method should be clearly distinguished from Multisource Holography, a system recently proposed by \citet{grace2023multi} for speckle reduction that also leverages a multi-source laser array. In Multisource Holography, the multi-source laser and two \textit{phase-only} SLMs placed in close proximity are used to remove speckles, but the system \'etendue remains limited since the spacing between each source is relatively small. In our system, we place the laser sources much farther apart to create a high-\'etendue backlight for the phase-only SLM to greatly increase the eyebox size and place an \textit{amplitude} display at the Fourier plane.